\chapter{Conclusions}
\label{cha:conclusion}

\section{Summary}
\label{sec:summary}

This thesis set out to determine whether frozen multimodal foundation models, fused deeply,
can forecast photovoltaic plants they have never observed. Answering that required building
three things: a dataset, a protocol, and a model.

The \textbf{dataset} fuses four public sources into $1{,}337{,}654$ rows over 110 sites in
two regions, pairing measured power and fourteen covariates with co-registered satellite
imagery on one timestamp grid, addressed by a single verified frame pointer, split by plant
identity with disjointness asserted at load time. The \textbf{protocol} fixes windows in
physical time, normalises by nameplate capacity so that no test-plant statistic enters any
fitted component, forbids domain physics inside models, and constrains retrieval datastores
to training plants. The \textbf{benchmark} runs 28 forecasters under that protocol, spanning
statistical references, classical machine learning, supervised deep networks, four
time-series foundation-model families, retrieval and adapter adaptation, and both endogenous
and exogenous multimodal systems. The \textbf{model}, MMTSFM, fuses frozen Chronos-2 and
frozen V-JEPA~2.1 through selective temporal interleaving and causal Grassmann mixing, under
a four-stage curriculum designed to keep randomly initialised modules from corrupting
pretrained residual streams.

\section{Answering the research question}
\label{sec:answer}

The research question asked whether deep token-level fusion of frozen foundation models
achieves cross-plant forecasting better than late fusion or domain-specific architectures.
The evidence currently available supports a partial answer, and the parts differ in how firm
they are.

\paragraph{H3 --- cross-plant transfer --- is established as a workable evaluation regime.}
Models trained on 69 plants forecast fourteen unseen plants from two weeks of history at
skill levels up to 0.552. The protocol discriminates sharply: the spread between the best and
worst models is far larger than any plausible measurement noise, so the benchmark is
informative rather than saturated.

\paragraph{H0 --- foundation models as anchors --- resolves negatively in its strong form.}
Zero-shot time-series foundation models do not merely underperform; three of four fall below
hourly climatology, and one falls below the naive floor. The result replicates across three
architecture families, so it is a property of photovoltaic dynamics rather than of a
particular pretraining corpus. Adaptation is mandatory on this task.

\paragraph{H4 --- retrieval --- is supported, and strongly.} Retrieval over a frozen backbone
reaches 0.478 against 0.331 for fine-tuning the same backbone: 0.147 skill, with no weight
updates. Among the frozen-backbone strategies tested, supplying relevant precedent beats both
re-fitting the model and injecting covariates through an adapter.

\paragraph{H1 and H2 --- vision, and fusion depth --- remain open, and the benchmark makes
them harder rather than easier.} The proposed model's own arms had not completed at the time
of writing, so the direct tests are pending. But the surrounding evidence is not neutral. The
best multimodal model in the suite consumes no external imagery, and every system that does
consume genuine satellite frames sits mid-pack. Any claim that deep fusion of real imagery
helps must now clear a bar set by a model that renders its own input as a picture.

\begin{table}[h!]
\centering
\small
\begin{tabular}{l p{6.6cm} p{4.6cm}}
\hline
\textbf{Hypothesis} & \textbf{Claim} & \textbf{Verdict} \\
\hline
H0 & A time-series foundation model applied zero-shot is a meaningful reference &
\textbf{Rejected in its strong form}: three of four fall below climatology \\
H1 & A visual stream improves over an identical time-series-only model & \textbf{Open};
the direct control was still running \\
H2 & Selective interleaving improves over late fusion & \textbf{Open}; same reason \\
H3 & Cross-plant zero-shot forecasting from a short history is achievable &
\textbf{Supported}: up to 0.552 skill on disjoint plants \\
H4 & Retrieval on a frozen backbone closes the gap to fine-tuning & \textbf{Supported, and
then some}: 0.478 against 0.331 \\
\hline
\end{tabular}
\caption{The hypothesis ladder of Section~\ref{sec:rq} against the evidence obtained.}
\label{tab:verdicts}
\end{table}

\section{What the benchmark teaches}
\label{sec:lessons}

Three conclusions generalise beyond this application.

\paragraph{Visual inductive bias and visual information are different resources, and the
literature conflates them.} Endogenous multimodal methods --- those that render a series as
an image --- exploit the \emph{priors} of a visual backbone: locality, periodicity, texture
sensitivity, multi-scale structure. Exogenous methods supply genuinely new observations. On
this benchmark the first converts into skill and the second, so far, does not. Reporting them
in one table, as is standard, obscures the finding entirely. At minimum, multimodal
forecasting results should state which resource is being exploited.

\paragraph{The modality-off arm is not optional.} A multimodal model that beats an external
unimodal baseline has demonstrated nothing about its second modality, because the comparison
confounds modality with training recipe, normalisation, horizon head and a dozen other
choices. Only an arm that is identical except for the modality isolates the contribution.
This thesis treats that arm as the primary result rather than as an appendix ablation, and
the design of Chapter~\ref{cha:expdesign} is built around it.

\paragraph{Conditioning may matter more than architecture.} The single largest measured
quantity in the entire study is not an architectural difference. It is the 0.173 skill gap
between a Chronos-2 variant that sees observed future weather and the same model that does
not. No pair of architectures in the 28-model suite differs by that much. If a fraction of
that gap can be recovered by better conditioning --- improved covariate forecasts, retrieved
analogous regimes, or a visual stream that actually informs the horizon --- it would exceed
what any architectural change in this study achieved.

\paragraph{Negative and null results carry the information here.} Three of the most useful
outputs of this work are things that did not happen: zero-shot foundation models did not
transfer, exogenous imagery did not convert into skill, and pretraining did not beat a
well-matched supervised architecture. None of these would have appeared in a study designed
only to demonstrate a proposed model, because none of them requires the proposed model to
observe. They required a benchmark with the alternatives in it. That is an argument for
building the instrument before building the method, and it is the methodological position
this thesis ended up defending.

\paragraph{A note on what ``foundation model'' bought.} It is worth separating two claims that
are often merged. Pretraining did not give zero-shot competence on this task --- that is
settled by Section~\ref{sec:finding2}. But pretraining did give a usable
\emph{initialisation}: fine-tuned Chronos-2 reaches 0.331 and adapted TTM-R2 0.364, both well
above their zero-shot counterparts and above the naive references. The honest summary is that
on photovoltaic data these models are good starting points and poor finished products, which
is a considerably weaker claim than the one the term ``foundation model'' invites.

\section{Limitations}
\label{sec:limitations}

\begin{itemize}
  \item \textbf{One dataset.} All results are \texttt{uk\_pv}: one region, one sensor, one
    capacity class, two years. Cross-dataset transfer is untested.
  \item \textbf{One seed per configuration}, no significance testing, and no measured
    seed-to-seed variance; small differences must be read as ties.
  \item \textbf{Half the control battery unexecuted}, including the two controls
    (shuffled frames, modality grid) that would most directly settle the vision question.
  \item \textbf{An open checkpoint-integrity defect} that blocks all post-hoc analysis of
    trained models (Section~\ref{sec:integrity}).
  \item \textbf{Random rather than region-based plant splits}, which allows spatially
    adjacent plants to straddle the split boundary.
  \item \textbf{An inconsistent \texttt{goes\_pvdaq} split} that predates the bad-site flags
    and must be regenerated before that track is used.
  \item \textbf{Two harness-limited baseline rows} whose numbers are indicative rather than
    precise.
\end{itemize}

\subsection{Which limitations actually threaten the conclusions}

The list above is not uniformly serious, and treating it as such would be its own form of
dishonesty. Ranked by how much they threaten what this thesis claims:

\emph{The single dataset is the real threat.} Every conclusion in
Chapter~\ref{cha:benchmark} is a conclusion about a fleet of 98 UK residential rooftops
observed through one satellite product over two years. The finding that supervised training
beats pretraining is exactly the kind of finding that could invert on a fleet with fewer
training plants or greater heterogeneity --- and the sample-efficiency scenario that would
test it was not run. This limitation bounds the scope of every headline in the thesis.

\emph{The missing significance testing is moderate.} It affects the reading of small
differences --- whether TS-RAG genuinely beats Cross-RAG --- but not the reading of large
ones. A 0.221 gap between iTransformer and fine-tuned Chronos-2 is not a statistical artefact
under any plausible variance estimate.

\emph{The checkpoint defect is severe but narrow.} It forecloses post-hoc analysis of trained
models entirely, which is a serious methodological handicap for the architecture work. It
does not affect any baseline number, since those were re-derived from saved predictions.

\emph{Random rather than region-based splits inflate all models roughly equally.} The concern
is real for the absolute values and weak for the ordering, which is what the thesis argues
from.

\emph{The unexecuted controls do not threaten a conclusion --- they withhold one.} Nothing
claimed here depends on them; they are what would be needed to claim more.

\section{Future work}
\label{sec:future}

The following are ordered by expected information gained per unit of compute, which is not
the order in which they are most tempting.

\begin{enumerate}
  \item \textbf{Run the two decisive controls.} Shuffled frames and the four-way modality
    grid are cheap relative to a training run and would determine whether the visual path is
    uninformative or merely mis-interfaced. Every subsequent decision about the vision branch
    depends on that answer, so it should be bought first.
  \item \textbf{Attack conditioning rather than architecture.} The 0.173 skill oracle gap is
    the largest measured headroom in this work. Concretely: better future-covariate
    estimates, regime retrieval as conditioning, or a visual stream whose output is injected
    at the horizon rather than only in the context.
  \item \textbf{Combine retrieval with fusion.} The two most effective levers identified ---
    retrieval over a frozen backbone, and multimodal conditioning --- have never been tested
    together. Retrieval supplies precedent, vision supplies the current state of the sky, and
    there is no obvious reason the two should be redundant.
  \item \textbf{Close the checkpoint-integrity defect} by activation-fingerprint bisection.
    Until this is fixed, no trained model in this line of work can be analysed after the fact,
    which forecloses interpretability work entirely.
  \item \textbf{Extend to the second dataset and to region-based splits.} Cross-dataset
    transfer (scenario S3) is the standard evidence a foundation-model claim is expected to
    provide, and region-based splits would remove the spatial-leakage concern in one step.
  \item \textbf{Reconsider the visual interface.} If the controls show that the frames carry
    signal the current encoder does not transmit, the summariser --- not the fusion operator
    --- is the component to redesign. A structured state formulation, in which the visual
    stream updates an explicit latent state rather than contributing tokens, is the natural
    next design and exists in draft.
\end{enumerate}

\subsection{The three experiments a follow-up should run first}

The list above is ordered by expected information per unit of compute. The top three deserve
concrete specification, because ``future work'' that is not specified does not get done.

\paragraph{Experiment 1: shuffled frames, one training run.} Take the interleaved arm, apply a
random temporal permutation to the visual window at evaluation only, and re-score. No
retraining is required, so the cost is one evaluation pass. Under the interpretation table of
Chapter~\ref{cha:expdesign}, an unchanged score is decisive: the model is not reading temporal
structure in the frames, and every downstream design decision about the visual branch changes.

\paragraph{Experiment 2: the modality grid, four training runs.} Target only; target plus
covariates; target plus vision; all three. Identical recipe, identical schedule, one seed each
--- ideally three. This is the experiment that attributes any gain to a specific input rather
than to their combination, and it subsumes several of the single zero-out ablations that would
otherwise be run separately. Four runs is the largest single item on this list and the one most
worth paying for.

\paragraph{Experiment 3: retrieval on top of the fusion model.} Add a TS-RAG-style retrieval
head over the frozen backbone of the interleaved arm, with the datastore restricted to
training plants. The two most effective levers identified in this work --- retrieval, worth
0.147 skill over fine-tuning, and multimodal conditioning --- have never been tested together,
and there is no obvious reason they should be redundant: retrieval supplies precedent, vision
supplies the current state of the sky. If they compose additively, the combination would clear
the 0.55 bar.

\subsection{If the visual interface is the problem}

Should the controls show that the frames carry signal the current encoder does not transmit,
the component to redesign is the summariser, not the fusion operator --- and three directions
are worth naming.

\emph{Widen the bottleneck before redesigning it.} A single 768-dimensional token stands in
for $4 \times 196$ patch embeddings (Table~\ref{tab:shapes}). Increasing the budget is the
cheapest possible test and should precede any architectural change.

\emph{Give the visual stream a task.} The summariser is currently trained only by the
forecasting loss, which is a weak and distant signal for learning what to extract from a
satellite crop. An auxiliary objective with a direct target --- predicting the next frame's
representation, or the observed clearness index at the site --- would supervise the interface
directly rather than through twelve encoder blocks and a quantile head.

\emph{Condition the horizon, not just the context.} Visual tokens currently enter only the
context region, so their influence on a horizon step is mediated entirely by the temporal
operator. Since cloud advection is a statement about the \emph{future} --- what is upwind now
will be overhead in two hours --- injecting a visual-derived representation at the future
positions is arguably the more natural design, and it is closer to how the oracle-covariate
variants achieve their 0.173 skill advantage. A structured state formulation, in which the
visual stream updates an explicit latent state carried forward across the horizon rather than
contributing context tokens, is the systematic version of this idea and exists in draft.

\section{Closing remark}
\label{sec:closing}

A thesis is expected to end with the model it proposed. This one ends somewhere else, and the
displacement is the honest summary of what was learned.

The most useful output of this work is not the proposed architecture but the instrument built
to evaluate it. A benchmark that ranks 28 forecasters under one protocol on disjoint plants
makes several comfortable assumptions in the multimodal forecasting literature testable, and
two of them do not survive contact: that pretrained foundation models transfer better than
supervised training on a well-matched architecture, and that consuming real imagery is what
makes a multimodal forecaster work. Whichever way the proposed model's own arms fall, those
two results stand, and they are what a subsequent effort should be designed around.

\par

